% 第5章 5.2.1：由半正定核构造 RKHS 的四个阶段（示意图，无数值）。
% 系数组合（半内积 a^T G b）-> 商去零长度方向得有限核和空间 H_0 -> 完备化 -> 点值读取 J 单射，得 H_K。
\begin{tikzpicture}[font=\figurefont\footnotesize,>={Stealth[length=2mm]},
  box/.style={draw=#1,fill=#1Fill,rounded corners=1.8pt,line width=.7pt,minimum width=21.5mm,minimum height=13mm,align=center,inner sep=1.2mm},
  note/.style={font=\FigureNoteFont\scriptsize,text=FigureMuted,align=center}]
\node[box=FigureBlue] (A) {系数组合\\[.6mm]$\displaystyle\sum_i a_iK_{x_i}$};
\node[box=FigureBlue,right=6mm of A] (B) {有限核和$H_0$\\[.6mm]真实函数};
\node[box=FigureGreen,right=6mm of B] (C) {完备化$\overline{H_0}$\\[.6mm]Cauchy列的类};
\node[box=FigureYellow,right=6mm of C] (D) {$H_K$\\[.6mm]$J\xi(x)=\langle\xi,K_x\rangle$};
\draw[->,FigureStroke,line width=.7pt] (A) -- (B);
\draw[->,FigureStroke,line width=.7pt] (B) -- (C);
\draw[->,FigureStroke,line width=.7pt] (C) -- (D);
\node[note,below=1.2mm of A] {半内积$\bm a^{\top}\bm G\bm b$\\可有零长度};
\node[note,below=1.2mm of B] {商去零长度\\得到内积};
\node[note,below=1.2mm of C] {点值随极限\\收敛};
\node[note,below=1.2mm of D] {核截面稠密\\$J$为单射};
\end{tikzpicture}
